\documentclass[11pt]{article}

\usepackage[margin=1in]{geometry}
\usepackage{amsmath}
\usepackage{booktabs}
\usepackage{array}
\usepackage{hyperref}

\title{Mathematical Description of the Unicorn \texttt{XJ} Tag}
\author{Julian Regalado Perez\\\texttt{julian.perez@sund.ku.dk}}
\date{}

\begin{document}
\maketitle

\section{Prompt}
\begin{quote}
\ttfamily\raggedright
Ok, I need you to provide me a latx file containing a thorough description of
the math behind the XJ tag alignment score recomputation.

Remember the implementaiton is found un
\texttt{unicorn\_internal.c::unicorn\_addjscore()}.

Add a small overall intro, an intro to the notation followed by method
description.

At the end add a working example with an alignment of a damaged fragment and
another one with an undamagged fragment.

Add the queries explicitly. Include both $5'$ C-T and $3'$ G-A. Place the
document in \texttt{XJtag.tex}.
\end{quote}

\section{Overview}

The Unicorn \texttt{XJ} tag is a damage-aware floating point alignment penalty
computed by the C function
\[
\texttt{unicorn\_addjscore()}
\]
in \path{src/libunicorn/unicorn_internal.c}. The purpose of the tag is to
recompute the mismatch component of an alignment score in a way that is less
punitive for substitutions that are characteristic of ancient DNA damage:
$5'$ C-to-T substitutions and $3'$ G-to-A substitutions.

The implementation should be read as a penalty score rather than a reward
score. Lower values indicate either fewer discrepancies or discrepancies that
are more compatible with terminal damage. Ordinary mismatches receive a penalty
of $1$, while damage-compatible mismatches receive a penalty between $0$ and
$1$ depending on their distance from the relevant fragment end. Insertions and
deletions are penalized with affine-like gap constants supplied by the caller.
In the current Unicorn callers, those constants are
\[
g_o = 5, \qquad g_e = 2,
\]
where $g_o$ is the gap opening penalty and $g_e$ is the per-base gap extension
penalty.

The final value is stored as a SAM/BAM auxiliary tag
\[
\texttt{XJ:f:<score>}.
\]
If an \texttt{XJ} tag already exists, it is deleted and replaced by the newly
computed value.

\section{Notation}

Consider one aligned BAM record. Let
\[
Q = Q_0 Q_1 \cdots Q_{L-1}
\]
be the query sequence stored in the BAM record, with query length $L$. Query
positions are zero-based, so $Q_0$ is the first base and $Q_{L-1}$ is the last
base. The implementation uses the query coordinates and query bases exactly as
provided by \texttt{bam\_get\_seq()} and \texttt{bam\_seqi()}.

The alignment is described by two pieces of SAM/BAM information:

\begin{itemize}
\item The CIGAR string, which describes how query columns and reference columns
are related.
\item The \texttt{MD} tag, which gives the reference bases at mismatch
positions and the reference bases deleted from the alignment.
\end{itemize}

Let $r$ denote a reference-column index in the alignment coordinate system used
internally by the function. The implementation builds an array
\[
\operatorname{ref2q}(r)
\]
that maps each reference-consuming alignment column to a query position:

\[
\operatorname{ref2q}(r) =
\begin{cases}
q, & \text{if reference column } r \text{ aligns to query base } Q_q,\\
-1, & \text{if reference column } r \text{ is a deletion from the query},\\
-2, & \text{if reference column } r \text{ is a reference skip}.
\end{cases}
\]

Only mismatch positions with $\operatorname{ref2q}(r) \ge 0$ contribute
mismatch penalties. Deleted reference bases are handled by the gap term, not by
the mismatch term.

The default damage model constants in the implementation are
\[
\alpha = 1.0, \qquad \lambda = 5.0,
\]
where $\alpha$ is the ordinary mismatch penalty and $\lambda$ controls the
exponential decay of the terminal damage discount.

\section{Method}

\subsection{Gap contribution}

The function first scans the CIGAR string. Insertions and deletions contribute
gap penalties. For an insertion of length $\ell$,
\[
S_I(\ell) = g_o + g_e \ell.
\]
For a deletion of length $\ell$,
\[
S_D(\ell) = g_o + g_e \ell.
\]

Thus the total gap contribution is
\[
S_{\mathrm{gap}}
 =
\sum_{I \in \mathcal{I}} \left(g_o + g_e \ell_I\right)
+
\sum_{D \in \mathcal{D}} \left(g_o + g_e \ell_D\right),
\]
where $\mathcal{I}$ is the set of CIGAR insertion operations and
$\mathcal{D}$ is the set of CIGAR deletion operations.

Soft clips advance the query coordinate but do not add a penalty. Reference
skips consume reference columns but do not add a penalty. Matches, sequence
matches, and sequence mismatches consume both reference and query columns but do
not directly add a gap penalty.

\subsection{Mismatch contribution}

Mismatch positions are recovered by parsing the \texttt{MD} tag. When the
\texttt{MD} parser sees a number, it advances by that many matching reference
columns. When it sees a deletion block beginning with \texttt{\^}, it advances
the reference coordinate through the deleted bases. When it sees an alphabetic
base, that base is the reference base at a mismatch position.

For a mismatch at reference column $r$, let
\[
q = \operatorname{ref2q}(r)
\]
be the aligned query coordinate. Let $R_r$ be the reference base from the
\texttt{MD} tag and let $Q_q$ be the query base from the BAM sequence. If
either base is \texttt{N}, the mismatch is ignored by the implementation.

The damage probability assigned to the mismatch is
\[
p_{\mathrm{damage}}(R_r, Q_q, q, L)
=
\begin{cases}
\exp\left(-q / \lambda\right),
& \text{if } R_r=\mathrm{C} \text{ and } Q_q=\mathrm{T},\\
\exp\left(-(L-q-1) / \lambda\right),
& \text{if } R_r=\mathrm{G} \text{ and } Q_q=\mathrm{A},\\
0,
& \text{otherwise.}
\end{cases}
\]

This means that C-to-T mismatches are discounted most strongly near the $5'$
end, while G-to-A mismatches are discounted most strongly near the $3'$ end.
The implementation clamps this probability to the interval $[0,1]$ after
computing it.

The mismatch penalty is then
\[
S_{\mathrm{mm}}(R_r, Q_q, q, L)
=
\alpha \left(1 - p_{\mathrm{damage}}(R_r, Q_q, q, L)\right).
\]
With $\alpha=1$, ordinary mismatches have penalty $1$. A terminal
damage-compatible mismatch can have penalty $0$, because
$\exp(0)=1$ at the relevant end.

\subsection{Final \texttt{XJ} value}

Let $\mathcal{M}$ be the set of mismatch positions reported by the \texttt{MD}
tag that map to real query bases and do not contain \texttt{N}. The final score
stored in the \texttt{XJ} tag is
\[
XJ
=
S_{\mathrm{gap}}
+
\sum_{r \in \mathcal{M}}
S_{\mathrm{mm}}\left(R_r, Q_{\operatorname{ref2q}(r)},
\operatorname{ref2q}(r), L\right).
\]

The function does not use the original aligner score, the \texttt{AS} tag, or
the edit distance tag \texttt{NM}. The \texttt{XJ} value is recomputed from the
CIGAR string, the query sequence, and the \texttt{MD} tag. If the BAM record is
missing an \texttt{MD} tag, no new \texttt{XJ} tag is written.

\section{Working Examples}

In both examples below, the CIGAR string is \texttt{16M}, so there are no
insertions or deletions and
\[
S_{\mathrm{gap}} = 0.
\]
The query length is $L=16$, the default damage decay length is $\lambda=5$, and
the ordinary mismatch penalty is $\alpha=1$.

\subsection{Damaged fragment with 5-prime C-to-T and 3-prime G-to-A}

Reference:
\[
\texttt{ACGTACCGTACGTAGC}
\]

Query:
\[
\texttt{ATGTACCGTACGTAAC}
\]

Alignment summary:

\begin{center}
\begin{tabular}{c c c c c}
\toprule
Position $q$ & Reference base & Query base & Damage class & Penalty \\
\midrule
1  & C & T & $5'$ C-to-T & $1-\exp(-1/5)$ \\
14 & G & A & $3'$ G-to-A & $1-\exp(-(16-14-1)/5)$ \\
\bottomrule
\end{tabular}
\end{center}

The corresponding \texttt{MD} tag is
\[
\texttt{MD:Z:1C12G1}.
\]
It means: one matching base, then a reference C mismatch, then twelve matching
bases, then a reference G mismatch, then one final matching base.

For the $5'$ C-to-T mismatch at $q=1$,
\[
S_{\mathrm{mm}} = 1 - \exp(-1/5)
               = 1 - \exp(-0.2)
               \approx 0.18127.
\]

For the $3'$ G-to-A mismatch at $q=14$, the distance from the $3'$ end is
\[
L-q-1 = 16 - 14 - 1 = 1,
\]
so
\[
S_{\mathrm{mm}} = 1 - \exp(-1/5)
               \approx 0.18127.
\]

Therefore
\[
XJ
= 0 + 0.18127 + 0.18127
\approx 0.36254.
\]

The resulting tag would be approximately
\[
\texttt{XJ:f:0.36254}.
\]

\subsection{Undamaged fragment with an ordinary internal mismatch}

Reference:
\[
\texttt{ACGTACCGTACGTAGC}
\]

Query:
\[
\texttt{ACGTGCCGTACGTAGC}
\]

This query has one internal A-to-G mismatch at query position $q=4$. It is not
a C-to-T mismatch and it is not a G-to-A mismatch, so it receives no damage
discount.

The corresponding \texttt{MD} tag is
\[
\texttt{MD:Z:4A11}.
\]
It means: four matching bases, then a reference A mismatch, then eleven
matching bases.

For this mismatch,
\[
p_{\mathrm{damage}} = 0,
\]
and therefore
\[
S_{\mathrm{mm}} = 1(1 - 0) = 1.
\]

Because there are no gaps,
\[
XJ = 0 + 1 = 1.
\]

The resulting tag would be
\[
\texttt{XJ:f:1}.
\]

If the undamaged query were an exact match to the reference,
\[
\texttt{ACGTACCGTACGTAGC},
\]
with \texttt{MD:Z:16}, then there would be no mismatch contribution and no gap
contribution, so $XJ=0$.

\end{document}
